\scanpage{58}
3: \quad Relations avec courbes elliptiques\inserted{~: le morphisme $f_2^t:\Sigma_{\underline J}\longrightarrow\Sigma_{\underline J}$}.

On suppose tjrs les car. rés. de $S\ne 2$. Soient $E_0$ une courbe elliptique sur $S$, $E$ un torseur sous $E_0$ (de sorte que $E_0\simeq\underline{\mathrm{Pic}}^0_{E/S}$).

On suppose donné dans $E$ un \inserted{$S$-}\unsure{autom|automorphisme}
\[
 \sigma:E\longrightarrow E \tag*{(17)}
\]
tel que
\[
 \sigma_{E_0}^{\unsure{2}}=-\mathrm{id}_{E_0} \tag*{(18)}
\]
ce qui implique que loc. (ét), identifiant $E$ à $E_0$ via une section origine de $E$, on aura
\[
 \sigma x=a-x\qquad a\ \text{section convenable}
\]
ce qui est d'ailleurs intrinsèque si on interprète $a$ comme une section de $E\mathbin{\wedge}_{E_0}E\simeq\underline{\mathrm{Pic}}^2_{E/S}$.

On aura donc
\[
 \sigma^2=\mathrm{id}_E
\]
Soit \deleted{$\unsure{\underline I|4}$}
\[
 \left\{\begin{aligned}
 &\underline I=E^\sigma\quad\text{sous-schéma des pts fixes de }\sigma\\
 &\underline I\subset E\\
 &\ \Vert\\[-3pt]
 &\{x\in E\mid 2x=a\}
 \end{aligned}\right. \tag*{(19)}
\]
on voit que $I$ est un torseur sous le schéma en groupes
\[
 \underline V={}_2E_0=\mathrm{Ker}\,(2\,\mathrm{id}_{E_0}:E_0\unsure{\text{—}|\longrightarrow} E_0), \tag*{(20)}
\]
qui est justement un groupe \inserted{sur $S$} loc. isom. à $\mathbb D_{2S}$.

La donnée de $(E,\sigma)$, où $E$ est un schéma
% archive sheet number: 57
